\documentclass[conference]{IEEEtran}
\IEEEoverridecommandlockouts
% The preceding line is only needed to identify funding in the first footnote. If that is unneeded, please comment it out.
%Template version as of 6/27/2024

\usepackage{cite}
\usepackage{amsmath,amssymb,amsfonts}
\usepackage{algorithmic}
\usepackage{graphicx}
\usepackage{textcomp}
\usepackage{xcolor}
\def\BibTeX{{\rm B\kern-.05em{\sc i\kern-.025em b}\kern-.08em
    T\kern-.1667em\lower.7ex\hbox{E}\kern-.125emX}}
\begin{document}

\title{Multi-View Convolutional Neural Network for Musculoskeletal Disorder Detection in Radiographic Imaging*\\
}

\author{\IEEEauthorblockN{1\textsuperscript{st} Camille Trouvé}
\IEEEauthorblockA{\textit{Ecole du Numérique (EDN)} \\
\textit{Faculty of Management, Economy and Science (FGES)}\\
Lille, France \\
trouve.camille@gmail.com}
\and
\IEEEauthorblockN{2\textsuperscript{nd} Ceren DINC}
\IEEEauthorblockA{\textit{Ecole du Numérique (EDN)} \\
\textit{Faculty of Management, Economy and Science (FGES)}\\
Lille, France \\
cdinc60@gmail.com}
\and
\IEEEauthorblockN{3\textsuperscript{rd} Esteban BOULANGER}
\IEEEauthorblockA{\textit{Ecole du Numérique (EDN)} \\
\textit{Faculty of Management, Economy and Science (FGES)}\\
Lille, France \\
es.boulanger@orange.fr}
\and
\IEEEauthorblockN{4\textsuperscript{th} Roger ACHKAR, Ph.D}
\IEEEauthorblockA{\textit{Faculty of Engineering & Technology} \\
\textit{Antonine University}\\
Baabda, Lebanon \\
dean.fet@ua.edu.lb }
}

\maketitle

\begin{abstract}
Musculoskeletal disorders are among the main causes of disability, according to the World Health Organization. This study aims to develop a tool capable of determining abnormalities in musculoskeletal radiographs. The proposed model is a convolutional neural network based on ResNet18 architecture with a multiple-view fusion module to integrate multiple radiographic projections, mimicking how radiologists interpret radiographic views. The model was trained and tested on the MURA dataset (Stanford), a large open-access dataset combining 40,561 radiographic images of the upper extremity from multiple studies. Through systematic evaluation of seven experimental configurations, the best-performing model achieved a validation AUC-ROC of 0.6545, representing a substantial improvement of +0.0906 compared to the baseline.
\end{abstract}

\begin{IEEEkeywords}
Musculoskeletal disorders, Radiography, Deep learning, CNN, Multi-view fusion
\end{IEEEkeywords}

\section{Introduction}
Musculoskeletal conditions represent a profound global health challenge, affecting approximately 1.7 billion people and serving as the primary driver of disability worldwide\cite{rajpurkar_mura_2018}. These disorders, which include fractures, osteoarthritis, and various degenerative joint diseases, significantly impair mobility and dexterity, often leading to premature exit from the workforce and a diminished quality of life. The increasing prevalence of these conditions, exacerbated by a growing and aging population, has placed an unprecedented strain on healthcare systems\cite{chada_machine_2019}. Radiologists  workloads rise because of a higher volume of complex images, which contributes to fatigue and a statistically significant risk of diagnostic error\cite{rajpurkar_mura_2018,chea_current_2020,singh_hybrid_2022}.

In this context, artificial intelligence (AI), specifically deep learning (DL) utilizing convolutional neural networks (CNNs), has emerged as a tool for diagnostic assistance\cite{chada_machine_2019,chea_current_2020,dangelo_artificial_2022}. CNN architectures are uniquely suited for image-based problems because they can autonomously learn discriminative features from large datasets, matching or even exceeding human performance in narrow tasks\cite{chea_current_2020,singh_hybrid_2022}. Within musculoskeletal radiology, DL applications have expanded across four primary categories: lesion detection, classification, segmentation, and non-interpretive tasks such as protocol optimization. While expert-level performance has been achieved in specific domains like bone age assessment and certain fracture detections, comprehensive interpretation of complex imaging examinations remains an ongoing challenge\cite{chea_current_2020,dangelo_artificial_2022}.

A pivotal milestone in this field was the release of the MURA (Musculoskeletal Radiographs) dataset, a large-scale collection of 40,561 upper extremity images manually labeled by board-certified radiologists\cite{rajpurkar_mura_2018}. Initial benchmark studies using a 169-layer DenseNet achieved high overall AUC-ROC scores but revealed a significant performance gap in specific anatomical regions, particularly the finger and humerus radiographs, where the model struggled to match the agreement levels of practicing radiologists. Furthermore, many existing models operate on single radiographic projections, whereas clinical diagnosis typically necessitates the integration of multiple views to resolve spatial ambiguities\cite{rajpurkar_mura_2018,dangelo_artificial_2022}.

This study proposes a multi-view neural network architecture designed to bridge this gap by mimicking the diagnostic process of human experts. By employing a ResNet18 backbone pretrained on ImageNet\cite{he_deep_2016} and a specialized masked average pooling module, the model integrates features from multiple radiographic views into a singular study-level prediction.


\section{Material and Method}

\item The subscript for the permeability of vacuum $\mu_{0}$, and other common scientific constants, is zero with subscript formatting, not a lowercase letter ``o''.
\item In American English, commas, semicolons, periods, question and exclamation marks are located within quotation marks only when a complete thought or name is cited, such as a title or full quotation. When quotation marks are used, instead of a bold or italic typeface, to highlight a word or phrase, punctuation should appear outside of the quotation marks. A parenthetical phrase or statement at the end of a sentence is punctuated outside of the closing parenthesis (like this). (A parenthetical sentence is punctuated within the parentheses.)
\item A graph within a graph is an ``inset'', not an ``insert''. The word alternatively is preferred to the word ``alternately'' (unless you really mean something that alternates)
\item Do not use the word ``essentially'' to mean ``approximately'' or ``effectively''.
\item In your paper title, if the words ``that uses'' can accurately replace the word ``using'', capitalize the ``u''; if not, keep using lower-cased.
\item Be aware of the different meanings of the homophones ``affect'' and ``effect'', ``complement'' and ``compliment'', ``discreet'' and ``discrete'', ``principal'' and ``principle''.
\item Do not confuse ``imply'' and ``infer''.
\item The prefix ``non'' is not a word; it should be joined to the word it modifies, usually without a hyphen.
\item There is no period after the ``et'' in the Latin abbreviation ``et al.''.
\item The abbreviation ``i.e.'' means ``that is'', and the abbreviation ``e.g.'' means ``for example''.
\end{itemize}
An excellent style manual for science writers is \cite{b7}.

\subsection{Authors and Affiliations}\label{AAA}
\textbf{The class file is designed for, but not limited to, six authors.} A 
minimum of one author is required for all conference articles. Author names 
should be listed starting from left to right and then moving down to the 
next line. This is the author sequence that will be used in future citations 
and by indexing services. Names should not be listed in columns nor group by 
affiliation. Please keep your affiliations as succinct as possible (for 
example, do not differentiate among departments of the same organization).

\subsection{Identify the Headings}\label{ITH}
Headings, or heads, are organizational devices that guide the reader through 
your paper. There are two types: component heads and text heads.

Component heads identify the different components of your paper and are not 
topically subordinate to each other. Examples include Acknowledgments and 
References and, for these, the correct style to use is ``Heading 5''. Use 
``figure caption'' for your Figure captions, and ``table head'' for your 
table title. Run-in heads, such as ``Abstract'', will require you to apply a 
style (in this case, italic) in addition to the style provided by the drop 
down menu to differentiate the head from the text.

Text heads organize the topics on a relational, hierarchical basis. For 
example, the paper title is the primary text head because all subsequent 
material relates and elaborates on this one topic. If there are two or more 
sub-topics, the next level head (uppercase Roman numerals) should be used 
and, conversely, if there are not at least two sub-topics, then no subheads 
should be introduced.

\subsection{Figures and Tables}\label{FAT}
\paragraph{Positioning Figures and Tables} Place figures and tables at the top and 
bottom of columns. Avoid placing them in the middle of columns. Large 
figures and tables may span across both columns. Figure captions should be 
below the figures; table heads should appear above the tables. Insert 
figures and tables after they are cited in the text. Use the abbreviation 
``Fig.~\ref{fig}'', even at the beginning of a sentence.

\begin{table}[htbp]
\caption{Table Type Styles}
\begin{center}
\begin{tabular}{|c|c|c|c|}
\hline
\textbf{Table}&\multicolumn{3}{|c|}{\textbf{Table Column Head}} \\
\cline{2-4} 
\textbf{Head} & \textbf{\textit{Table column subhead}}& \textbf{\textit{Subhead}}& \textbf{\textit{Subhead}} \\
\hline
copy& More table copy$^{\mathrm{a}}$& &  \\
\hline
\multicolumn{4}{l}{$^{\mathrm{a}}$Sample of a Table footnote.}
\end{tabular}
\label{tab1}
\end{center}
\end{table}

\begin{figure}[htbp]
\centerline{\includegraphics{fig1.png}}
\caption{Example of a figure caption.}
\label{fig}
\end{figure}

Figure Labels: Use 8 point Times New Roman for Figure labels. Use words 
rather than symbols or abbreviations when writing Figure axis labels to 
avoid confusing the reader. As an example, write the quantity 
``Magnetization'', or ``Magnetization, M'', not just ``M''. If including 
units in the label, present them within parentheses. Do not label axes only 
with units. In the example, write ``Magnetization (A/m)'' or ``Magnetization 
\{A[m(1)]\}'', not just ``A/m''. Do not label axes with a ratio of 
quantities and units. For example, write ``Temperature (K)'', not 
``Temperature/K''.

\section*{Acknowledgment}

The preferred spelling of the word ``acknowledgment'' in America is without 
an ``e'' after the ``g''. Avoid the stilted expression ``one of us (R. B. 
G.) thanks $\ldots$''. Instead, try ``R. B. G. thanks$\ldots$''. Put sponsor 
acknowledgments in the unnumbered footnote on the first page.

\section*{References}

Please number citations consecutively within brackets \cite{b1}. The 
sentence punctuation follows the bracket \cite{b2}. Refer simply to the reference 
number, as in \cite{b3}---do not use ``Ref. \cite{b3}'' or ``reference \cite{b3}'' except at 
the beginning of a sentence: ``Reference \cite{b3} was the first $\ldots$''

Number footnotes separately in superscripts. Place the actual footnote at 
the bottom of the column in which it was cited. Do not put footnotes in the 
abstract or reference list. Use letters for table footnotes.

Unless there are six authors or more give all authors' names; do not use 
``et al.''. Papers that have not been published, even if they have been 
submitted for publication, should be cited as ``unpublished'' \cite{b4}. Papers 
that have been accepted for publication should be cited as ``in press'' \cite{b5}. 
Capitalize only the first word in a paper title, except for proper nouns and 
element symbols.

For papers published in translation journals, please give the English 
citation first, followed by the original foreign-language citation \cite{b6}.

\begin{thebibliography}{00}
\bibitem{b1} G. Eason, B. Noble, and I. N. Sneddon, ``On certain integrals of Lipschitz-Hankel type involving products of Bessel functions,'' Phil. Trans. Roy. Soc. London, vol. A247, pp. 529--551, April 1955.
\bibitem{b2} J. Clerk Maxwell, A Treatise on Electricity and Magnetism, 3rd ed., vol. 2. Oxford: Clarendon, 1892, pp.68--73.
\bibitem{b3} I. S. Jacobs and C. P. Bean, ``Fine particles, thin films and exchange anisotropy,'' in Magnetism, vol. III, G. T. Rado and H. Suhl, Eds. New York: Academic, 1963, pp. 271--350.
\bibitem{b4} K. Elissa, ``Title of paper if known,'' unpublished.
\bibitem{b5} R. Nicole, ``Title of paper with only first word capitalized,'' J. Name Stand. Abbrev., in press.
\bibitem{b6} Y. Yorozu, M. Hirano, K. Oka, and Y. Tagawa, ``Electron spectroscopy studies on magneto-optical media and plastic substrate interface,'' IEEE Transl. J. Magn. Japan, vol. 2, pp. 740--741, August 1987 [Digests 9th Annual Conf. Magnetics Japan, p. 301, 1982].
\bibitem{b7} M. Young, The Technical Writer's Handbook. Mill Valley, CA: University Science, 1989.
\bibitem{b8} D. P. Kingma and M. Welling, ``Auto-encoding variational Bayes,'' 2013, arXiv:1312.6114. [Online]. Available: https://arxiv.org/abs/1312.6114
\bibitem{b9} S. Liu, ``Wi-Fi Energy Detection Testbed (12MTC),'' 2023, gitHub repository. [Online]. Available: https://github.com/liustone99/Wi-Fi-Energy-Detection-Testbed-12MTC
\bibitem{b10} ``Treatment episode data set: discharges (TEDS-D): concatenated, 2006 to 2009.'' U.S. Department of Health and Human Services, Substance Abuse and Mental Health Services Administration, Office of Applied Studies, August, 2013, DOI:10.3886/ICPSR30122.v2
\bibitem{b11} K. Eves and J. Valasek, ``Adaptive control for singularly perturbed systems examples,'' Code Ocean, Aug. 2023. [Online]. Available: https://codeocean.com/capsule/4989235/tree
\end{thebibliography}

\vspace{12pt}
\color{red}
IEEE conference templates contain guidance text for composing and formatting conference papers. Please ensure that all template text is removed from your conference paper prior to submission to the conference. Failure to remove the template text from your paper may result in your paper not being published.

\end{document}
